\documentclass[11pt]{article}
\usepackage[a4paper,margin=0.8in]{geometry}
\usepackage{calc}
\usepackage{amsmath,amssymb}
\usepackage{array,longtable,booktabs}
\usepackage{xcolor}
\usepackage{hyperref}
\usepackage{enumitem}
\usepackage{fancyvrb}
\usepackage{fvextra}
\usepackage{framed}
\usepackage{color}
\fvset{breaklines=true, breakanywhere=true}
\DefineVerbatimEnvironment{Highlighting}{Verbatim}{breaklines,breakanywhere,commandchars=\\\{\}}
\usepackage{color}
\usepackage{fancyvrb}
\newcommand{\VerbBar}{|}
\newcommand{\VERB}{\Verb[commandchars=\\\{\}]}
\DefineVerbatimEnvironment{Highlighting}{Verbatim}{commandchars=\\\{\}}
% Add ',fontsize=\small' for more characters per line
\usepackage{framed}
\definecolor{shadecolor}{RGB}{248,248,248}
\newenvironment{Shaded}{\begin{snugshade}}{\end{snugshade}}
\newcommand{\AlertTok}[1]{\textcolor[rgb]{0.94,0.16,0.16}{#1}}
\newcommand{\AnnotationTok}[1]{\textcolor[rgb]{0.56,0.35,0.01}{\textbf{\textit{#1}}}}
\newcommand{\AttributeTok}[1]{\textcolor[rgb]{0.13,0.29,0.53}{#1}}
\newcommand{\BaseNTok}[1]{\textcolor[rgb]{0.00,0.00,0.81}{#1}}
\newcommand{\BuiltInTok}[1]{#1}
\newcommand{\CharTok}[1]{\textcolor[rgb]{0.31,0.60,0.02}{#1}}
\newcommand{\CommentTok}[1]{\textcolor[rgb]{0.56,0.35,0.01}{\textit{#1}}}
\newcommand{\CommentVarTok}[1]{\textcolor[rgb]{0.56,0.35,0.01}{\textbf{\textit{#1}}}}
\newcommand{\ConstantTok}[1]{\textcolor[rgb]{0.56,0.35,0.01}{#1}}
\newcommand{\ControlFlowTok}[1]{\textcolor[rgb]{0.13,0.29,0.53}{\textbf{#1}}}
\newcommand{\DataTypeTok}[1]{\textcolor[rgb]{0.13,0.29,0.53}{#1}}
\newcommand{\DecValTok}[1]{\textcolor[rgb]{0.00,0.00,0.81}{#1}}
\newcommand{\DocumentationTok}[1]{\textcolor[rgb]{0.56,0.35,0.01}{\textbf{\textit{#1}}}}
\newcommand{\ErrorTok}[1]{\textcolor[rgb]{0.64,0.00,0.00}{\textbf{#1}}}
\newcommand{\ExtensionTok}[1]{#1}
\newcommand{\FloatTok}[1]{\textcolor[rgb]{0.00,0.00,0.81}{#1}}
\newcommand{\FunctionTok}[1]{\textcolor[rgb]{0.13,0.29,0.53}{\textbf{#1}}}
\newcommand{\ImportTok}[1]{#1}
\newcommand{\InformationTok}[1]{\textcolor[rgb]{0.56,0.35,0.01}{\textbf{\textit{#1}}}}
\newcommand{\KeywordTok}[1]{\textcolor[rgb]{0.13,0.29,0.53}{\textbf{#1}}}
\newcommand{\NormalTok}[1]{#1}
\newcommand{\OperatorTok}[1]{\textcolor[rgb]{0.81,0.36,0.00}{\textbf{#1}}}
\newcommand{\OtherTok}[1]{\textcolor[rgb]{0.56,0.35,0.01}{#1}}
\newcommand{\PreprocessorTok}[1]{\textcolor[rgb]{0.56,0.35,0.01}{\textit{#1}}}
\newcommand{\RegionMarkerTok}[1]{#1}
\newcommand{\SpecialCharTok}[1]{\textcolor[rgb]{0.81,0.36,0.00}{\textbf{#1}}}
\newcommand{\SpecialStringTok}[1]{\textcolor[rgb]{0.31,0.60,0.02}{#1}}
\newcommand{\StringTok}[1]{\textcolor[rgb]{0.31,0.60,0.02}{#1}}
\newcommand{\VariableTok}[1]{\textcolor[rgb]{0.00,0.00,0.00}{#1}}
\newcommand{\VerbatimStringTok}[1]{\textcolor[rgb]{0.31,0.60,0.02}{#1}}
\newcommand{\WarningTok}[1]{\textcolor[rgb]{0.56,0.35,0.01}{\textbf{\textit{#1}}}}
\setlist[itemize]{leftmargin=*}
\setlist[enumerate]{leftmargin=*}
\hypersetup{colorlinks=true,linkcolor=blue,urlcolor=blue}
\providecommand{\tightlist}{%
  \setlength{\itemsep}{0pt}\setlength{\parskip}{0pt}}
\title{R Workshop -- Day 6 -- Coding Statistical Models from Their Equations}
\author{Applied R for Statistical Teaching, Research and Reproducible Analysis}
\date{}
\begin{document}
\maketitle
\tableofcontents
\newpage
\hypertarget{r-workshop-day-6-coding-statistical-models-from-their-equations}{%
\section{R Workshop -- Day 6 -- Coding Statistical Models from Their
Equations}\label{r-workshop-day-6-coding-statistical-models-from-their-equations}}

\textbf{Session 6 of 7 \textbar{} Duration: 3 hours} \textbf{Programme:}
Applied R for Statistical Teaching, Research and Reproducible Analysis

\hypertarget{learning-objectives}{%
\subsection{Learning objectives}\label{learning-objectives}}

By the end of this session you will be able to:

\begin{enumerate}
\def\labelenumi{\arabic{enumi}.}
\tightlist
\item
  Read ordinary least squares (OLS) notation and identify the dimensions
  of each matrix involved.
\item
  Build a design matrix from a formula and a data frame.
\item
  Derive and compute OLS coefficient estimates directly from matrix
  algebra.
\item
  Calculate fitted values and residuals from first principles.
\item
  Implement a small, reusable custom OLS function.
\item
  Validate that function against \texttt{lm()}.
\item
  Discuss one numerical failure mode, such as rank deficiency or poor
  conditioning.
\end{enumerate}

\hypertarget{capstone-link}{%
\subsection{Capstone link}\label{capstone-link}}

Today you open the ``black box'' behind the Day 5 model: you reproduce
(or partially audit) its coefficients by computing them directly from
the design matrix, rather than from \texttt{lm()}.

\hypertarget{segment-plan-3-hours}{%
\subsection{Segment plan (3 hours)}\label{segment-plan-3-hours}}

\begin{longtable}[]{@{}
  >{\raggedright\arraybackslash}p{(\columnwidth - 4\tabcolsep) * \real{0.3000}}
  >{\raggedleft\arraybackslash}p{(\columnwidth - 4\tabcolsep) * \real{0.4000}}
  >{\raggedright\arraybackslash}p{(\columnwidth - 4\tabcolsep) * \real{0.3000}}@{}}
\toprule\noalign{}
\begin{minipage}[b]{\linewidth}\raggedright
Segment
\end{minipage} & \begin{minipage}[b]{\linewidth}\raggedleft
Time
\end{minipage} & \begin{minipage}[b]{\linewidth}\raggedright
Purpose
\end{minipage} \\
\midrule\noalign{}
\endhead
\bottomrule\noalign{}
\endlastfoot
Opening and recap & 10 min & Recap the Day 5 model, preview today's
equation-to-code work \\
Concept bridge & 25 min & OLS notation, dimensions, the normal
equations \\
Guided coding & 60 min & Design matrix, manual OLS, custom function \\
Participant practice & 45 min & Validate against \texttt{lm()}, test one
failure mode \\
Interpretation and reporting & 25 min & Explain the equation in your own
words \\
Evidence log and capstone update & 15 min & Save the equation-based
implementation \\
\end{longtable}

\begin{center}\rule{0.5\linewidth}{0.5pt}\end{center}

\hypertarget{part-1-concept-bridge-core-workshop-activity}{%
\subsection{Part 1 -- Concept bridge (Core workshop
activity)}\label{part-1-concept-bridge-core-workshop-activity}}

\hypertarget{notation-and-dimensions}{%
\subsubsection{Notation and dimensions}\label{notation-and-dimensions}}

For \(n\) observations and \(p\) predictors (including the intercept),
ordinary least squares assumes:

\[y = X\beta + \varepsilon\]

where \(y\) is an \(n \times 1\) vector of outcomes, \(X\) is an
\(n \times p\) \textbf{design matrix} (a column of 1s for the intercept,
plus one column per predictor), \(\beta\) is a \(p \times 1\) vector of
coefficients, and \(\varepsilon\) is an \(n \times 1\) vector of errors.

\hypertarget{the-normal-equations}{%
\subsubsection{The normal equations}\label{the-normal-equations}}

OLS chooses \(\hat\beta\) to minimise the sum of squared residuals,
\((y - X\beta)^\top(y - X\beta)\). Differentiating and setting to zero
gives the \textbf{normal equations}:

\[X^\top X \hat\beta = X^\top y \quad\Longrightarrow\quad \hat\beta = (X^\top X)^{-1} X^\top y\]

This single line is the algebraic content of every \texttt{lm()} call
you have run this week.

\textbf{Important implementation note:} the formula above assumes \(X\)
includes an intercept column and is full rank (no exact linear
dependence among predictor columns). Our first custom function will
assume an intercept explicitly; a model fitted without one needs
different degrees-of-freedom and R-squared bookkeeping, which is out of
scope today (see Reference material).

\begin{center}\rule{0.5\linewidth}{0.5pt}\end{center}

\hypertarget{part-2-guided-coding-core-workshop-activity}{%
\subsection{Part 2 -- Guided coding (Core workshop
activity)}\label{part-2-guided-coding-core-workshop-activity}}

\begin{Shaded}
\begin{Highlighting}[]
\NormalTok{lecturers }\OtherTok{\textless{}{-}}\NormalTok{ readr}\SpecialCharTok{::}\FunctionTok{read\_csv}\NormalTok{(}\StringTok{"Data/processed/anchor\_dataset\_clean.csv"}\NormalTok{, }\AttributeTok{show\_col\_types =} \ConstantTok{FALSE}\NormalTok{)}
\NormalTok{model\_data }\OtherTok{\textless{}{-}}\NormalTok{ lecturers[stats}\SpecialCharTok{::}\FunctionTok{complete.cases}\NormalTok{(}
\NormalTok{  lecturers[, }\FunctionTok{c}\NormalTok{(}\StringTok{"post\_score"}\NormalTok{, }\StringTok{"pre\_score"}\NormalTok{, }\StringTok{"attendance"}\NormalTok{)]}
\NormalTok{), ]}
\end{Highlighting}
\end{Shaded}

We use the simpler two-predictor model
\texttt{post\_score\ \textasciitilde{}\ pre\_score\ +\ attendance}
today, so the matrix algebra stays easy to follow; the same method
extends to the full Day 5 model.

\hypertarget{build-the-design-matrix}{%
\subsubsection{2.1 Build the design
matrix}\label{build-the-design-matrix}}

\begin{Shaded}
\begin{Highlighting}[]
\NormalTok{X }\OtherTok{\textless{}{-}} \FunctionTok{model.matrix}\NormalTok{(}\SpecialCharTok{\textasciitilde{}}\NormalTok{ pre\_score }\SpecialCharTok{+}\NormalTok{ attendance, }\AttributeTok{data =}\NormalTok{ model\_data)}
\NormalTok{y }\OtherTok{\textless{}{-}}\NormalTok{ model\_data}\SpecialCharTok{$}\NormalTok{post\_score}

\FunctionTok{dim}\NormalTok{(X)     }\CommentTok{\# n rows, 3 columns: intercept, pre\_score, attendance}
\FunctionTok{head}\NormalTok{(X, }\DecValTok{3}\NormalTok{)}
\end{Highlighting}
\end{Shaded}

\texttt{model.matrix()} does exactly what you would otherwise do by
hand: add an intercept column of 1s and bind the predictor columns
alongside it.

\hypertarget{compute-ols-coefficients-directly}{%
\subsubsection{2.2 Compute OLS coefficients
directly}\label{compute-ols-coefficients-directly}}

\begin{Shaded}
\begin{Highlighting}[]
\NormalTok{XtX }\OtherTok{\textless{}{-}} \FunctionTok{t}\NormalTok{(X) }\SpecialCharTok{\%*\%}\NormalTok{ X}
\NormalTok{XtY }\OtherTok{\textless{}{-}} \FunctionTok{t}\NormalTok{(X) }\SpecialCharTok{\%*\%}\NormalTok{ y}
\NormalTok{beta\_hat }\OtherTok{\textless{}{-}} \FunctionTok{solve}\NormalTok{(XtX) }\SpecialCharTok{\%*\%}\NormalTok{ XtY}
\NormalTok{beta\_hat}
\end{Highlighting}
\end{Shaded}

Compare directly:

\begin{Shaded}
\begin{Highlighting}[]
\NormalTok{model\_lm }\OtherTok{\textless{}{-}} \FunctionTok{lm}\NormalTok{(post\_score }\SpecialCharTok{\textasciitilde{}}\NormalTok{ pre\_score }\SpecialCharTok{+}\NormalTok{ attendance, }\AttributeTok{data =}\NormalTok{ model\_data)}
\FunctionTok{coef}\NormalTok{(model\_lm)}
\end{Highlighting}
\end{Shaded}

The two sets of numbers should match to several decimal places.

\hypertarget{fitted-values-and-residuals}{%
\subsubsection{2.3 Fitted values and
residuals}\label{fitted-values-and-residuals}}

\begin{Shaded}
\begin{Highlighting}[]
\NormalTok{fitted\_manual }\OtherTok{\textless{}{-}}\NormalTok{ X }\SpecialCharTok{\%*\%}\NormalTok{ beta\_hat}
\NormalTok{residuals\_manual }\OtherTok{\textless{}{-}}\NormalTok{ y }\SpecialCharTok{{-}}\NormalTok{ fitted\_manual}

\CommentTok{\# Compare with lm()}
\FunctionTok{max}\NormalTok{(}\FunctionTok{abs}\NormalTok{(fitted\_manual }\SpecialCharTok{{-}} \FunctionTok{fitted}\NormalTok{(model\_lm)))}
\FunctionTok{max}\NormalTok{(}\FunctionTok{abs}\NormalTok{(residuals\_manual }\SpecialCharTok{{-}} \FunctionTok{residuals}\NormalTok{(model\_lm)))}
\end{Highlighting}
\end{Shaded}

\hypertarget{a-small-custom-ols-function}{%
\subsubsection{2.4 A small custom OLS
function}\label{a-small-custom-ols-function}}

\begin{Shaded}
\begin{Highlighting}[]
\NormalTok{my\_lm }\OtherTok{\textless{}{-}} \ControlFlowTok{function}\NormalTok{(formula, data) \{}
\NormalTok{  X }\OtherTok{\textless{}{-}} \FunctionTok{model.matrix}\NormalTok{(formula, }\AttributeTok{data =}\NormalTok{ data)}
\NormalTok{  response\_name }\OtherTok{\textless{}{-}} \FunctionTok{all.vars}\NormalTok{(formula)[}\DecValTok{1}\NormalTok{]}
\NormalTok{  y }\OtherTok{\textless{}{-}}\NormalTok{ data[[response\_name]][stats}\SpecialCharTok{::}\FunctionTok{complete.cases}\NormalTok{(}\FunctionTok{model.matrix}\NormalTok{(formula, }\AttributeTok{data =}\NormalTok{ data))]}

  \ControlFlowTok{if}\NormalTok{ (}\FunctionTok{colnames}\NormalTok{(X)[}\DecValTok{1}\NormalTok{] }\SpecialCharTok{!=} \StringTok{"(Intercept)"}\NormalTok{) \{}
    \FunctionTok{stop}\NormalTok{(}\StringTok{"my\_lm() requires an intercept; interceptless formulas are not supported."}\NormalTok{)}
\NormalTok{  \}}

\NormalTok{  beta\_hat }\OtherTok{\textless{}{-}} \FunctionTok{solve}\NormalTok{(}\FunctionTok{t}\NormalTok{(X) }\SpecialCharTok{\%*\%}\NormalTok{ X) }\SpecialCharTok{\%*\%}\NormalTok{ (}\FunctionTok{t}\NormalTok{(X) }\SpecialCharTok{\%*\%}\NormalTok{ y)}
\NormalTok{  fitted\_values }\OtherTok{\textless{}{-}}\NormalTok{ X }\SpecialCharTok{\%*\%}\NormalTok{ beta\_hat}
\NormalTok{  resid\_values }\OtherTok{\textless{}{-}}\NormalTok{ y }\SpecialCharTok{{-}}\NormalTok{ fitted\_values}

\NormalTok{  n }\OtherTok{\textless{}{-}} \FunctionTok{nrow}\NormalTok{(X)}
\NormalTok{  p }\OtherTok{\textless{}{-}} \FunctionTok{ncol}\NormalTok{(X)}
\NormalTok{  rss }\OtherTok{\textless{}{-}} \FunctionTok{sum}\NormalTok{(resid\_values}\SpecialCharTok{\^{}}\DecValTok{2}\NormalTok{)}
\NormalTok{  tss }\OtherTok{\textless{}{-}} \FunctionTok{sum}\NormalTok{((y }\SpecialCharTok{{-}} \FunctionTok{mean}\NormalTok{(y))}\SpecialCharTok{\^{}}\DecValTok{2}\NormalTok{)}
\NormalTok{  r\_squared }\OtherTok{\textless{}{-}} \DecValTok{1} \SpecialCharTok{{-}}\NormalTok{ rss }\SpecialCharTok{/}\NormalTok{ tss}
\NormalTok{  sigma2 }\OtherTok{\textless{}{-}}\NormalTok{ rss }\SpecialCharTok{/}\NormalTok{ (n }\SpecialCharTok{{-}}\NormalTok{ p)}
\NormalTok{  var\_beta }\OtherTok{\textless{}{-}}\NormalTok{ sigma2 }\SpecialCharTok{*} \FunctionTok{solve}\NormalTok{(}\FunctionTok{t}\NormalTok{(X) }\SpecialCharTok{\%*\%}\NormalTok{ X)}
\NormalTok{  se\_beta }\OtherTok{\textless{}{-}} \FunctionTok{sqrt}\NormalTok{(}\FunctionTok{diag}\NormalTok{(var\_beta))}

  \FunctionTok{list}\NormalTok{(}
    \AttributeTok{coefficients =} \FunctionTok{setNames}\NormalTok{(}\FunctionTok{as.vector}\NormalTok{(beta\_hat), }\FunctionTok{colnames}\NormalTok{(X)),}
    \AttributeTok{std\_errors =} \FunctionTok{setNames}\NormalTok{(se\_beta, }\FunctionTok{colnames}\NormalTok{(X)),}
    \AttributeTok{fitted\_values =} \FunctionTok{as.vector}\NormalTok{(fitted\_values),}
    \AttributeTok{residuals =} \FunctionTok{as.vector}\NormalTok{(resid\_values),}
    \AttributeTok{r\_squared =}\NormalTok{ r\_squared,}
    \AttributeTok{df\_residual =}\NormalTok{ n }\SpecialCharTok{{-}}\NormalTok{ p}
\NormalTok{  )}
\NormalTok{\}}
\end{Highlighting}
\end{Shaded}

\hypertarget{validate-against-lm}{%
\subsubsection{\texorpdfstring{2.5 Validate against
\texttt{lm()}}{2.5 Validate against lm()}}\label{validate-against-lm}}

\begin{Shaded}
\begin{Highlighting}[]
\NormalTok{my\_fit }\OtherTok{\textless{}{-}} \FunctionTok{my\_lm}\NormalTok{(post\_score }\SpecialCharTok{\textasciitilde{}}\NormalTok{ pre\_score }\SpecialCharTok{+}\NormalTok{ attendance, }\AttributeTok{data =}\NormalTok{ model\_data)}
\NormalTok{my\_fit}\SpecialCharTok{$}\NormalTok{coefficients}
\FunctionTok{coef}\NormalTok{(model\_lm)}

\NormalTok{my\_fit}\SpecialCharTok{$}\NormalTok{r\_squared}
\FunctionTok{summary}\NormalTok{(model\_lm)}\SpecialCharTok{$}\NormalTok{r.squared}

\NormalTok{my\_fit}\SpecialCharTok{$}\NormalTok{std\_errors}
\FunctionTok{summary}\NormalTok{(model\_lm)}\SpecialCharTok{$}\NormalTok{coefficients[, }\StringTok{"Std. Error"}\NormalTok{]}
\end{Highlighting}
\end{Shaded}

All three comparisons should agree closely. Small differences in the
last few decimal places are expected and come from differing numerical
routines, not from a mistake.

\begin{center}\rule{0.5\linewidth}{0.5pt}\end{center}

\hypertarget{part-3-participant-practice-core-workshop-activity}{%
\subsection{Part 3 -- Participant practice (Core workshop
activity)}\label{part-3-participant-practice-core-workshop-activity}}

\begin{enumerate}
\def\labelenumi{\arabic{enumi}.}
\item
  Run \texttt{my\_lm()} on the Day 5 model formula
  (\texttt{post\_score\ \textasciitilde{}\ pre\_score\ +\ attendance\ +\ training\_track})
  and confirm coefficients, R-squared and standard errors match
  \texttt{lm()}.

  \emph{Hint:} \texttt{training\_track} must be converted to numeric
  dummy columns -- \texttt{model.matrix()} does this automatically once
  \texttt{training\_track} is a factor.
\item
  Deliberately create a \textbf{rank-deficient} design matrix
  (e.g.~include both \texttt{attendance} and an exact duplicate column
  \texttt{attendance\_copy\ \textless{}-\ attendance}) and run it
  through \texttt{my\_lm()}. Read and explain the resulting error from
  \texttt{solve()}.
\item
  Investigate \textbf{poor conditioning}: compute
  \texttt{kappa(t(X)\ \%*\%\ X)} for the well-specified model and for a
  version where one predictor is rescaled to a tiny range
  (e.g.~\texttt{attendance\ /\ 10000}). Note how the condition number
  changes.
\item
  Write two or three sentences, in your own words, explaining what
  \((X^\top X)^{-1} X^\top y\) is doing -- as if explaining it to a
  colleague who knows regression results but has never seen the matrix
  form.
\end{enumerate}

A facilitator-reviewed solution is in
\texttt{Solution-Scripts/day6\_manual\_ols\_solution.R}.

\begin{center}\rule{0.5\linewidth}{0.5pt}\end{center}

\hypertarget{optional-demonstration}{%
\subsection{Optional demonstration}\label{optional-demonstration}}

\begin{itemize}
\tightlist
\item
  \textbf{Bootstrap uncertainty} -- resample rows with replacement,
  refit \texttt{my\_lm()} each time, and compare the bootstrap standard
  errors to the analytic ones above.
\item
  \textbf{Maximum likelihood} -- show that, under normal errors, the OLS
  estimator is also the maximum-likelihood estimator (a one-paragraph
  derivation, not a full implementation).
\item
  \textbf{Detailed conditioning analysis} --
  \texttt{solve(t(X)\ \%*\%\ X)} versus \texttt{MASS::ginv()} when
  \(X^\top X\) is nearly singular.
\item
  \textbf{Extended inference} -- compute \(t\)-statistics and p-values
  from \texttt{std\_errors} inside \texttt{my\_lm()} and compare to
  \texttt{summary(model\_lm)}.
\end{itemize}

\begin{Shaded}
\begin{Highlighting}[]
\CommentTok{\# Optional: bootstrap standard error for the attendance coefficient}
\FunctionTok{set.seed}\NormalTok{(}\DecValTok{2024}\NormalTok{)}
\NormalTok{boot\_coefs }\OtherTok{\textless{}{-}} \FunctionTok{replicate}\NormalTok{(}\DecValTok{500}\NormalTok{, \{}
\NormalTok{  boot\_rows }\OtherTok{\textless{}{-}} \FunctionTok{sample}\NormalTok{(}\FunctionTok{seq\_len}\NormalTok{(}\FunctionTok{nrow}\NormalTok{(model\_data)), }\AttributeTok{replace =} \ConstantTok{TRUE}\NormalTok{)}
\NormalTok{  boot\_fit }\OtherTok{\textless{}{-}} \FunctionTok{my\_lm}\NormalTok{(post\_score }\SpecialCharTok{\textasciitilde{}}\NormalTok{ pre\_score }\SpecialCharTok{+}\NormalTok{ attendance, }\AttributeTok{data =}\NormalTok{ model\_data[boot\_rows, ])}
\NormalTok{  boot\_fit}\SpecialCharTok{$}\NormalTok{coefficients[}\StringTok{"attendance"}\NormalTok{]}
\NormalTok{\})}
\FunctionTok{sd}\NormalTok{(boot\_coefs)}
\end{Highlighting}
\end{Shaded}

\begin{center}\rule{0.5\linewidth}{0.5pt}\end{center}

\hypertarget{reference-or-take-home-material}{%
\subsection{Reference or take-home
material}\label{reference-or-take-home-material}}

\begin{itemize}
\tightlist
\item
  \textbf{QR-based fitting:} \texttt{lm()} does not actually invert
  \(X^\top X\) -- it uses a QR decomposition of \(X\) for better
  numerical stability. Explicit matrix inversion (as we did today, for
  teaching clarity) is more sensitive to poor conditioning than QR-based
  methods, and is avoided in production statistical software.
\item
  \textbf{Interceptless-model extensions:} if a model genuinely has no
  intercept (\texttt{y\ \textasciitilde{}\ x\ -\ 1}), R-squared, degrees
  of freedom and the overall F-test all require different formulas; see
  \texttt{?summary.lm} for how R itself handles this case.
\end{itemize}

\begin{center}\rule{0.5\linewidth}{0.5pt}\end{center}

\hypertarget{daily-deliverable}{%
\subsection{Daily deliverable}\label{daily-deliverable}}

An equation-based implementation (\texttt{my\_lm()} and its outputs)
that reproduces or audits part of the Day 5 model, with a short written
explanation of the matrix algebra in your own words.

\hypertarget{evidence-log-entry}{%
\subsubsection{Evidence log entry}\label{evidence-log-entry}}

Complete the Day 6 row in \texttt{Workshop-Evidence-Log.md}: confirm
your custom function matches \texttt{lm()} and note what the
rank-deficiency exercise taught you about when a model cannot be fitted
at all.

\begin{center}\rule{0.5\linewidth}{0.5pt}\end{center}

\hypertarget{facilitator-notes}{%
\subsection{Facilitator notes}\label{facilitator-notes}}

\begin{itemize}
\tightlist
\item
  The rank-deficiency exercise should produce a clear, interpretable
  error from \texttt{solve()} (a singular-matrix error); if a
  participant's \texttt{attendance\_copy} is not an exact duplicate, the
  matrix may instead be merely ill-conditioned rather than singular --
  worth distinguishing live.
\item
  Numerical agreement with \texttt{lm()} should be exact to at least 6-8
  significant figures for the well-conditioned cases; larger
  discrepancies usually indicate a data-preparation mismatch
  (e.g.~different rows being dropped due to missingness) rather than an
  algebra error.
\item
  This session is mathematically the densest of the workshop. If the
  group needs more guided-coding time, drop the bootstrap optional
  demonstration first.
\end{itemize}
\end{document}
